% mage-auto-heading: up-front-encode-what-you-already-know
% mage-auto-heading: as-you-learn-convert-what-the-work-teaches
% mage-auto-heading: correct-judgment-without-enforcement
% mage-auto-heading: engineering-capital
% mage-auto-heading: where-enforcement-stops
% Authoritative LaTeX manuscript; Markdown is a generated compatibility projection.

% <!-- part-title: Alignment -->

% <!-- chapter-title: Growing the Governed Environment -->

\chapter{Growing the Governed Environment}

% <!-- point: a-governed-environment-grows-from-known-obligations-and-from-what-the-work-teaches | A governed environment grows from obligations already known and from what the work teaches; governance conversion is the second route. | terms: up-front-governance, governance-conversion -->

% <!-- index-def: governed-environment -->

No engineered environment begins with every future obligation already represented and mechanized. Some obligations are known in advance: an API must require authentication, a state transition is forbidden, a format mutation must cross a particular seam. Others become visible only when the work exposes a failure nobody represented or controlled. A governed environment acquires durable structure in two ways: \textbf{up front, from obligations already known, and in response to failures and surprises.}

% <!-- index-def: governance-conversion -->

The second route is \textbf{governance conversion}. A failure exposes a recurring class, engineering identifies what was missing, and the environment encodes the lesson as durable structure. Later work inherits the result. The repair fixes this instance. Conversion changes the environment.

% <!-- box-family: inset -->

% <!-- inset-domain: ENGINEERING -->

% mage-box-prefix: Inset —

\begin{bookbox}{inset}{Failure Modes Before and After Failure}
\label{inset-failure-modes}
Engineering need not wait for a failure to learn how a system can fail. Techniques such as Failure Modes and Effects Analysis (FMEA) work up front: engineers identify possible failure modes, consider their effects, and use that analysis to decide where prevention, detection, mitigation, or other controls are warranted. Known obligations can therefore produce durable engineering structure before experience forces the issue.

Governance conversion works in the complementary direction. A failure or important surprise supplies evidence that something was missing. Engineering identifies the failure mode behind the instance and asks what durable structure should address the class: a model, obligation, sensor, validator, constraint, gate, or other mechanism.

The distinction is about how the need became known, not about the resulting mechanism. Up-front analysis and learning from experience may lead to the same control. Governance conversion names the second route; it does not imply that engineers should wait for failures they can reasonably anticipate.
\end{bookbox}

\section{Up Front: Encode What You Already Know}\label{heading:up-front-encode-what-you-already-know}

% <!-- section-terms: up-front-governance, specification-driven-development -->

% <!-- point: up-front-encode-obligations-you-can-state-before-the-work-begins | Up front: encode the obligations you can state before the work begins, and stop paying judgment for them. | terms: up-front-governance, specification-driven-development -->

% <!-- index-def: up-front-governance -->

% <!-- index-def: specification-driven-development -->

Start with obligations you can state before the work begins and that are consequential, recurrent, or reusable enough to make durable. A required output shape may become a schema; a forbidden dependency may become an architectural relation plus a validator; a service boundary may become a permission or narrow API; a legal lifecycle may become a closed transition table. Encode recurring obligations when the environment can carry them more cheaply and reliably than repeated judgment.

\magevalue{chapter:one-problem-many-models} turns this into a method for starting and migrating systems; here the distinction is simply that these obligations are known before failure forces the issue.

\section{As You Learn: Convert What the Work Teaches}\label{heading:as-you-learn-convert-what-the-work-teaches}

% <!-- section-terms: governance-conversion, constraints-and-sensors -->

% <!-- point: diagnose-what-was-missing-do-not-assume-every-failure-wants-a-gate | Diagnose what was missing; not every failure calls for another gate. | terms: governance-conversion -->

The harder half begins when the system surprises you. \textbf{Do not assume every failure calls for another gate. Diagnose what was missing.}

% <!-- point: the-conversion-lands-in-one-of-five-places | The conversion lands in one of five places: model, invariant, sensor, validator, or enforcement. | terms: governance-conversion, constraints-and-sensors -->

The conversion can land in one of five places:

\begin{enumerate}
\def\labelenumi{\arabic{enumi}.}
\tightlist
\item
  \textbf{Missing representation / semantics → improve the model.} Make the fact or relation durable at the abstraction where the question becomes answerable.
\end{enumerate}

\begin{enumerate}
\def\labelenumi{\arabic{enumi}.}
\setcounter{enumi}{1}
\tightlist
\item
  \textbf{Missing obligation or tolerance → state the governing boundary.} A failure may show that variation treated as free was not actually acceptable. Make the invariant, policy, bound, or permitted set explicit enough for future work to reason against it.
\end{enumerate}

\begin{enumerate}
\def\labelenumi{\arabic{enumi}.}
\setcounter{enumi}{2}
\tightlist
\item
  \textbf{Missing evidence → add a sensor or trace.} Make the relevant state observable.
\end{enumerate}

\begin{enumerate}
\def\labelenumi{\arabic{enumi}.}
\setcounter{enumi}{3}
\tightlist
\item
  \textbf{Missing evaluation → strengthen the validator.} Turn evidence into a repeatable judgment.
\end{enumerate}

\begin{enumerate}
\def\labelenumi{\arabic{enumi}.}
\setcounter{enumi}{4}
\tightlist
\item
  \textbf{Missing enforcement → add or strengthen a constraint or gate.} Make the verdict control the relevant action or admission.
\end{enumerate}

These diagnoses are not mutually exclusive. In practice, one failure often exposes several missing pieces at once. A property that was never represented may also lack an explicit invariant and the sensor needed to observe it. The durable repair can therefore add a model, an obligation, and evidence together. The categories identify where engineering structure was missing; they do not require the repair to land in only one place.

These repairs address the problem introduced in \magevalue{chapter:problem}: how much consequential work must depend on an agent reconstructing the right knowledge and making the right judgment each time? A model can preserve distinctions the agent would otherwise have to infer; an independent validator can evaluate a recurring obligation without relying on the agent\textquotesingle s own assessment. Governance conversion makes such settled engineering decisions part of the environment that subsequent work inherits. Engineers can then concentrate judgment on new questions and on decisions whose assumptions have changed.

% <!-- point: modeling-and-alignment-form-a-feedback-loop | Agent work supplies evidence for engineering decisions about models, obligations, and controls; later work uses the revised environment and supplies further evidence. | terms: modeling-principle, alignment-principle, governance-conversion -->

Agents use the models and operate under the controls of this environment. Their results and observed behavior provide evidence for engineers to assess how well the environment supports the work. A failure or surprise may reveal knowledge that remained implicit, including a choice treated as a degree of freedom that actually carried a tacit obligation or tolerance. Engineers investigate the cause and decide whether the model, the stated obligation, or the control should change. Later work uses the revised environment and supplies further evidence, completing the feedback loop. \autoref{governance-conversion-fig} shows how diagnosis can lead to durable structure that benefits future work.

\begin{figure}[htbp]
\centering
\includegraphics[width=\linewidth]{book/latex/assets/governance-conversion-dbc51d550d.pdf}
\caption{\emph{Governance conversion.} When a failure exposes a recurring failure mode rooted in missing representation, obligation, evidence, evaluation, or enforcement, encode the missing structure at the appropriate layer. Future work inherits the result as engineering capital when that durable structure continues to lower future cost or uncertainty.}
\label{governance-conversion-fig}
\end{figure}

\section{Correct Judgment Without Enforcement}\label{heading:correct-judgment-without-enforcement}

% <!-- point: correct-judgment-without-enforcement | A correct validation result does not protect the system unless an appropriate mechanism acts on it. | terms: governance-conversion, validator, gate -->

Suppose a deployment pipeline checks that every required configuration value is present. The validator detects a missing value, but the pipeline records only a warning and continues deploying. The evidence and evaluation are adequate for that obligation; acceptance does not depend on the verdict. Adding another check for the same missing value would repeat the judgment without repairing the gap.

The durable change is to make the deployment decision depend on that result, if the configuration really is required. In another case, the safer repair may remove an unsafe operation entirely and require review when it would have been used. Governance conversion should change the part of the environment that permitted the failure. It can add a mechanism, strengthen its consequence, or retire a capability whose risk no longer has a useful justification.

\section{Engineering Capital}\label{heading:engineering-capital}

% <!-- section-terms: engineering-capital, governance-conversion, modeling-principle, alignment-principle -->

% <!-- point: governance-conversion-forms-engineering-capital-the-positive-stock | Governance conversion forms engineering capital — the concrete positive stock software engineering has lacked a name for. | terms: engineering-capital, governance-conversion -->

% <!-- index-def: engineering-capital -->

Governance conversion motivates a useful concept: \textbf{engineering capital}. You may be familiar with the opposite idea, technical debt: a shortcut today that makes future changes more costly. Where technical debt reduces your ability to make changes, engineering capital increases it. You have durable models, checks, and other engineering structure that lower the cost or uncertainty of future changes.

% <!-- point: the-asset-is-something-you-can-point-at | The capital must be something future work can actually inherit — a concrete asset you can point at. | terms: engineering-capital, constraints-and-sensors -->

% <!-- point: capital-earns-a-return-lower-marginal-engineering-cost | Capital earns a return: lower marginal engineering cost. Conversion is capital formation. | terms: engineering-capital, governance-conversion -->

The capital must be something future work can actually inherit. A model can eliminate repeated reconstruction; a constraint can prevent a recurring failure; a sensor can expose previously hidden state; a validator can make repeated evaluation mechanical; a gate can enforce that evaluation; a skill or runbook can package judgment for reuse. These assets change the starting conditions of future work. They earn a return when future work requires less reconstruction, repeated judgment, failure, or rework because the structure exists.

% <!-- point: engineering-capital-can-transfer-between-projects | Engineering capital can transfer through reusable mechanisms or through the design judgment behind them; its value depends on work saved after adaptation and maintenance costs. | terms: engineering-capital, modeling-principle, alignment-principle -->

Engineering capital can also transfer between projects. A reusable build check or task-completion procedure can give a new project an established starting point. Engineers may adapt the mechanism directly, or reuse the design judgment that led to it while building a different mechanism. The value comes from work the next project no longer has to repeat, after accounting for the cost of adapting and maintaining what it inherits.

% <!-- point: inherited-capital-improves-starting-conditions-without-supplying-every-control | Inherited automation and principles improve a new project's starting conditions; they do not supply every representation or control the new system needs. | terms: engineering-capital, governance-conversion -->

Reuse still requires engineering. A deployment check may assume a particular hosting platform; a review procedure may assume a team structure the new project lacks. The receiving project must identify those assumptions, test whether they hold, and adapt or discard the parts that do not fit. Inherited structure improves the starting conditions only where its obligations and mechanisms remain appropriate.

% <!-- point: capital-depreciates-so-accumulation-cannot-be-the-objective | Capital depreciates; maintenance, reconciliation, and retirement manage it, so accumulation cannot be the objective. | terms: engineering-capital, governed-environment -->

The engineered environment evolves with the system it governs.\magefootnote{lehman}{Lehman\textquotesingle s work on software evolution made the analogous observation about software itself: software used in a changing real-world environment must continually adapt, and its complexity tends to increase unless effort is spent maintaining or reducing it. MAGE extends that caution to the engineered environment surrounding the software. Models, validators, constraints, and procedures evolve too --- and can become liabilities when they no longer fit what they govern. See Lehman \& Ramil 2003.\citep{lehman2003evolution}} Capital also \textbf{depreciates}. A model can drift from the system it represents. A validator can encode an assumption the architecture has outgrown. A sensor can become noisy; a gate can impose more delay than the failure it prevents; a once-important control can survive long after its obligation disappears. Maintenance preserves useful capital; reconciliation repairs assets that no longer match the system; retirement removes assets whose carrying cost now exceeds their return. This is why accumulation cannot be the objective. A governed environment with more mechanisms is not necessarily richer than one with fewer. What matters is the future engineering value those mechanisms continue to produce.

% <!-- point: debt-and-capital-are-symmetric-stocks-that-compound | Technical debt and engineering capital are symmetric stocks — both accumulate, both change the economics of later change, both compound. | terms: engineering-capital -->

Debt makes future work pay again for today\textquotesingle s expedience; capital lets future work inherit today\textquotesingle s judgment. Governance conversion turns lessons from individual engineering efforts into structure that subsequent work can reuse.

\section{Where Enforcement Stops}\label{heading:where-enforcement-stops}

% <!-- section-terms: residual, constraints-and-sensors -->

% <!-- point: where-a-goal-lands-turns-on-three-questions-starting-with-explicit-in-usable-form | Where a goal lands turns on three questions; the first is whether the obligation is explicit in a form the environment can use. | terms: residual, constraints-and-sensors -->

Not every quality goal should end in enforcement. A consequential concern may remain unrepresented, require expert judgment, be evaluated without enforcement, or be governed. \autoref{honest-gaps} shows the four possible stopping points.

\begin{figure}[htbp]
\centering
\includegraphics[width=\linewidth]{book/latex/assets/residual-decision-tree-478c6933ab.pdf}
\caption{\emph{Where enforcement stops.} A concern may remain unrepresented, require judgment, be evaluated without enforcement, or be governed. These are design choices; different obligations should stop at different points.}
\label{honest-gaps}
\end{figure}

% <!-- point: the-residual-is-an-obligation-nobody-made-explicit-enough-for-the-environment-to-use | The residual is an obligation nobody made explicit enough for the environment to use. | terms: residual, judgment-is-the-scarce-resource -->

% <!-- index-def: residual -->

\textbf{Residual.} The obligation is not explicit enough for the environment to use, so an engineer or agent must reconstruct the missing meaning.

% <!-- point: an-explicit-obligation-may-still-resist-cheap-mechanical-evaluation | An explicit obligation may still resist cheap mechanical evaluation — an honest boundary, not failed Alignment. | terms: residual, constraints-and-sensors -->

\textbf{Judgment required.} The obligation is explicit, but available mechanisms cannot evaluate it adequately. Architecture may be represented while a design trade-off still requires expert judgment.

% <!-- point: observed-only-is-explicit-and-evaluated-but-un-gated-because-a-gate-is-a-commitment | Evaluated only: explicit and evaluated, yet un-gated — because a gate is a commitment. | terms: residual, constraints-and-sensors -->

\textbf{Evaluated only.} The environment can produce and evaluate relevant evidence, but does not enforce the resulting judgment.

% <!-- point: a-governed-obligation-carries-a-binding-consequence-appropriate-to-its-risk | A governed obligation is enforced in a form appropriate to the engineering decision. | terms: residual, constraints-and-sensors -->

\textbf{Governed.} The environment enforces the obligation in a form appropriate to the engineering decision: an action may be constrained, a violation refused, or admission made dependent on the verdict.

% <!-- point: stronger-alignment-is-not-narrower-tolerance | Making a tolerance explicit does not mean tightening it; stronger Alignment means dependable enforcement of the obligation engineering chose, not a narrower acceptable set. | terms: residual, tolerance -->

The same restraint applies to tolerance. Making an acceptable boundary explicit does not imply that the environment should tighten it, or even enforce it mechanically. A broad tolerance may be the correct engineering decision, and variation inside it remains legitimate realization freedom. Stronger Alignment means more dependable enforcement of the obligation engineering actually chose, not progressively narrowing the acceptable realization space.

% <!-- point: three-different-reasons-non-enforcement-is-sometimes-the-engineered-result | An obligation can stop short of a gate for three different reasons; non-enforcement is sometimes the engineered result. | terms: residual, governance-conversion -->

Different concerns can warrant different stopping points. A team may record service cost while leaving the spending decision to a human, rather than imposing a hard daily ceiling. It may review a trend in incident reports without setting a required number, which would reward the count instead of the improvement. A design decision may require context that the available checks cannot evaluate adequately. In each case, engineers should make the remaining judgment or policy choice explicit instead of treating non-enforcement as an accidental omission.

% <!-- point: the-mechanism-that-discovers-need-not-carry-the-enforcement | The mechanism that discovers an obligation need not be the one that ultimately enforces it; Chapter 4 follows that migration. | terms: sensor, governance-conversion -->

The mechanism that exposes a problem need not be the one that ultimately enforces the repair. A complexity warning may identify code worth examining without justifying a ban on every similar function. Investigation may instead reveal a missing type or an architectural boundary that a more precise control can enforce. \magevalue{chapter:one-problem-many-models} follows a migration of this kind, from an audit signal to more specific representations and controls.

% <!-- point: the-point-is-not-the-rightmost-leaf-it-is-knowing-where-each-obligation-sits-and-why | The point is not to drive every obligation to the rightmost leaf; it is to know where each sits and why — and conversion has now created an estate to govern. | terms: governance-conversion, residual, governed-environment -->

Different obligations should stop in different places. Governance conversion closes recurring gaps when doing so is worth the cost. It also creates the next problem: the environment now contains an expanding portfolio of controls, and those controls can themselves become inconsistent, redundant, or orphaned.
